% tuto_mpm_solide.tex -- Tutoriel : un solide MPM avec endommagement et rupture, couplé au fluide APIC
% Compiler avec pdfLaTeX (Overleaf : compilateur par défaut).
\documentclass[11pt,a4paper]{article}

\usepackage[utf8]{inputenc}
\usepackage[T1]{fontenc}
\usepackage{lmodern}
\usepackage[french]{babel}
\usepackage[margin=2.2cm]{geometry}
\usepackage{amsmath,amssymb}
\usepackage{xcolor}
\usepackage{booktabs}
\usepackage{tikz}
\usetikzlibrary{arrows.meta,positioning,calc}
\usepackage{listings}
\usepackage{hyperref}
\hypersetup{colorlinks=true,linkcolor=blue!60!black,urlcolor=blue!60!black}

% ---------------------------------------------------------------- style des listings
\definecolor{codebg}{RGB}{247,247,250}
\definecolor{codekw}{RGB}{0,80,160}
\definecolor{codecomment}{RGB}{90,120,90}
\definecolor{codestring}{RGB}{160,60,20}
\lstset{
  language=Python,
  basicstyle=\ttfamily\small,
  keywordstyle=\color{codekw}\bfseries,
  commentstyle=\color{codecomment}\itshape,
  stringstyle=\color{codestring},
  backgroundcolor=\color{codebg},
  frame=single, framerule=0pt, rulecolor=\color{codebg},
  xleftmargin=6pt, xrightmargin=6pt,
  breaklines=true, breakatwhitespace=true,
  showstringspaces=false, tabsize=4,
  columns=fullflexible, keepspaces=true,
  literate=
    {é}{{\'e}}1 {è}{{\`e}}1 {ê}{{\^e}}1 {ë}{{\"e}}1
    {à}{{\`a}}1 {â}{{\^a}}1 {î}{{\^i}}1 {ï}{{\"i}}1
    {ô}{{\^o}}1 {ù}{{\`u}}1 {û}{{\^u}}1 {ç}{{\c{c}}}1
    {É}{{\'E}}1 {È}{{\`E}}1 {À}{{\`A}}1 {œ}{{\oe}}1
}
\newcommand{\code}[1]{\texttt{\detokenize{#1}}}
\newcommand{\bm}[1]{\boldsymbol{#1}}
\newcommand{\tr}{\operatorname{tr}}

\title{Un solide MPM avec endommagement et rupture,\\ couplé au fluide APIC\\
  \large Tutoriel en deux parties : le pont seul, puis le pont sous l'eau}
\author{}
\date{}

\begin{document}
\maketitle

\begin{abstract}
Ce tutoriel ajoute au solveur APIC 2D du tutoriel précédent une \textbf{phase solide} capable de
se déformer, de s'endommager et de se rompre : un pont, un tablier, un mur. La méthode est
\textbf{MLS-MPM} (\emph{Material Point Method}), qui partage avec APIC la même grille, les mêmes
poids d'interpolation et le même cycle P2G $\to$ grille $\to$ G2P. Ce qui change tient en trois
ingrédients portés par chaque particule solide : un \textbf{gradient de déformation} $\mathbf F$
et une loi élastique, une \textbf{variable d'endommagement} $D\in[0,1]$ qui réduit la raideur,
et un état \textbf{rompu} qui supprime la cohésion.
La partie~I construit et teste le solide \emph{seul} (une poutre encastrée que l'on charge
jusqu'à la ruine), la partie~II le couple au fluide sur une grille commune (un bloc d'eau qui
tombe sur le tablier). Le code complet, testé avec Taichi 1.7.4, est donné en annexe et vit dans
le dossier \code{MpM/} du projet.
\end{abstract}

\tableofcontents
\bigskip

% ================================================================
\section{Prérequis et organisation des fichiers}

On suppose acquis le tutoriel APIC (poids B-spline quadratiques, \code{base}/\code{fx}/\code{w},
transfert affine $C_p$, formule $D_p^{-1} = 4/\Delta x^2$). Le projet est organisé ainsi :

\begin{center}
\begin{tabular}{@{}ll@{}}
\toprule
Fichier & Rôle \\
\midrule
\code{APIC/APIC.py} & fluide : \code{P2G}, \code{grid_step}, \code{G2P} (inchangé) \\
\code{MpM/mpm_solid.py} & solide : \code{P2G_solid}, \code{grid_update}, \code{G2P_solid}, utilitaires \\
\code{MpM/demo_solid.py} & partie I : la poutre seule, avec curseur de charge \\
\code{main_fsi.py} & partie II : eau + pont sur la grille commune \\
\bottomrule
\end{tabular}
\end{center}

Lancer depuis la racine du projet :
\begin{lstlisting}[language=bash]
python -m MpM.demo_solid     # partie I  : la poutre seule
python main_fsi.py           # partie II : eau + pont
\end{lstlisting}
Touches communes : \textbf{Espace} bascule la couleur des particules solides (endommagement /
déformation), \textbf{R} réinitialise, \textbf{Échap} quitte. Dans la partie I, les touches
\textbf{1}, \textbf{2}, \textbf{3} choisissent le modèle (élastique, $+$endommagement, $+$rupture)
et un curseur règle la charge.

% ================================================================
\section{L'idée générale}

\subsection{MPM, c'est APIC avec une mémoire de la déformation}

Une particule de fluide APIC transporte $(x_p, v_p, C_p, J_p)$ : sa seule mémoire de l'histoire
est le scalaire $J_p$, taux de dilatation, qui suffit à définir une pression. Un solide, lui,
résiste à \emph{toute} déformation, pas seulement au changement de volume : cisaillement, flexion,
étirement. Il faut donc une mémoire tensorielle, le \textbf{gradient de déformation} $\mathbf F_p$
($2\times2$), et une loi qui transforme $\mathbf F_p$ en contrainte $\bm\sigma_p$ ($2\times2$).
Le reste du cycle est identique :

\begin{center}
\begin{tikzpicture}[
  box/.style={draw, rounded corners, minimum width=4cm, minimum height=1.7cm, align=center},
  arr/.style={-{Latex[length=3mm]}, thick},
  lab/.style={font=\small, inner sep=2pt, fill=white}]
  \node[box, fill=blue!5] (f) at (0,1.8) {Particules fluide\\ $x_p,\,v_p,\,C_p,\,J_p$};
  \node[box, fill=orange!10] (s) at (0,-1.8) {Particules solide\\ $x_p,\,v_p,\,C_p$\\ $\mathbf F_p,\,D_p,\,\texttt{broken}_p$};
  \node[box, fill=gray!10, minimum height=5.2cm] (g) at (9.2,0)
    {Grille commune\\[3pt] $m_i = m_i^{f} + m_i^{s}$\\ $(mv)_i \to v_i$\\[3pt] gravité, parois,\\ piliers};
  \draw[arr] ($(f.east)+(0,0.4)$) -- node[lab, above]{P2G : $m$, $mv$, pression} ($(g.west|-f.east)+(0,0.4)$);
  \draw[arr, dashed] ($(g.west|-f.east)+(0,-0.4)$) -- node[lab, below]{G2P : $v$, $C$, $J$} ($(f.east)+(0,-0.4)$);
  \draw[arr] ($(s.east)+(0,0.4)$) -- node[lab, above]{P2G : $m$, $mv$, force élastique} ($(g.west|-s.east)+(0,0.4)$);
  \draw[arr, dashed] ($(g.west|-s.east)+(0,-0.4)$) -- node[lab, below]{G2P : $v$, $C$, $\mathbf F$, $D$, rupture} ($(s.east)+(0,-0.4)$);
\end{tikzpicture}
\end{center}

\subsection{Trois niveaux de modèle}

\begin{enumerate}
  \item \textbf{Élastique} : $\mathbf F_p$ est mis à jour à chaque pas, une loi hyperélastique donne la
        contrainte, la force interne est déposée sur la grille. La poutre fléchit, vibre, revient.
  \item \textbf{Endommagement} : une variable $D_p$ croît de façon irréversible quand l'allongement
        dépasse un seuil ; la raideur devient $(1-D_p)E$. La poutre s'assouplit là où elle est
        trop sollicitée.
  \item \textbf{Rupture} : quand $D_p = 1$ la particule perd toute cohésion : elle ne transmet plus de
        traction, seulement de la compression. Les morceaux se séparent et tombent.
\end{enumerate}
On implémente les trois d'un coup dans \code{G2P_solid} avec deux drapeaux \code{use_damage} et
\code{use_rupture}, pour pouvoir les comparer au clavier.

\subsection{Par rapport au code MPM classique}

Le code MPM \og éléments Q4 \fg{} que vous avez peut-être sous les yeux (fonctions de forme de
Lagrange, \code{nodal_iforces -= Vp * stress @ dNdx}, contrainte incrémentée par
\code{dsigma = C @ dEps}) fait la même chose que nous, avec trois différences de forme :

\begin{center}
\begin{tabular}{@{}lll@{}}
\toprule
 & MPM classique (Q4) & MLS-MPM (ce tutoriel) \\
\midrule
Fonctions de forme & Lagrange bilinéaires, $2\times2$ n\oe uds & B-spline quadratique, $3\times3$ n\oe uds \\
Gradient des poids & $\partial N_I/\partial x$ calculé via le jacobien & $\nabla w_{ip} \approx \frac{4}{\Delta x^2} w_{ip}(x_i - x_p)$ \\
Gradient de vitesse & $L_p = \sum_I v_I\, \nabla N_I^{\mathsf T}$ & $C_p$ (déjà calculé pour APIC) \\
Contrainte & hypoélastique : $\Delta\bm\sigma = \mathbf C\,\Delta\bm\varepsilon$ & hyperélastique : $\bm\sigma = \bm\sigma(\mathbf F)$ \\
\bottomrule
\end{tabular}
\end{center}

La formulation hyperélastique (contrainte fonction de $\mathbf F$ seul) est plus robuste aux
grandes rotations, ce qui compte pour des morceaux de pont qui tombent en tournant. Et
$C_p \approx \nabla v$ est \emph{gratuit} : APIC le calcule déjà.

% ================================================================
\section{Le minimum de mathématiques}

\subsection{Le gradient de déformation}
\label{sec:F}

Une particule solide occupait la position $\mathbf X_p$ au repos et occupe $\mathbf x_p$ maintenant.
Le gradient de déformation $\mathbf F = \partial\mathbf x/\partial\mathbf X$ décrit comment un petit
voisinage a été étiré et tourné : $\mathbf F = \mathbf I$ au repos, $\det\mathbf F = J$ est le rapport
de volume (le $J_p$ du fluide !). Son évolution suit
\[
  \dot{\mathbf F} = (\nabla v)\,\mathbf F,
\]
et comme $C_p \approx \nabla v$ au voisinage de la particule, la mise à jour discrète est simplement
\[
  \boxed{\;\mathbf F_p \leftarrow (\mathbf I + \Delta t\, C_p)\,\mathbf F_p\;}
\]
C'est la généralisation matricielle de $J_p \leftarrow J_p(1 + \Delta t\,\tr C_p)$ du fluide
(prendre le déterminant des deux côtés redonne cette formule au premier ordre).

\subsection{Décomposition polaire et allongements principaux}

Toute matrice $\mathbf F$ se décompose en une rotation et un étirement pur : $\mathbf F = \mathbf R\,\mathbf S$
avec $\mathbf R$ orthogonale et $\mathbf S$ symétrique. Numériquement on passe par la SVD
$\mathbf F = \mathbf U\,\bm\Sigma\,\mathbf V^{\mathsf T}$, d'où
\[
  \mathbf R = \mathbf U\mathbf V^{\mathsf T},\qquad
  \bm\Sigma = \operatorname{diag}(s_1, s_2).
\]
Les $s_k$ sont les \textbf{allongements principaux} : $s_k = 1$ au repos, $s_k > 1$ en traction,
$s_k < 1$ en compression, dans deux directions orthogonales. Taichi fournit \code{ti.svd(F)} pour
les matrices $2\times2$ et $3\times3$ ; on l'utilise partout.

\subsection{La loi de comportement}
\label{sec:loi}

On utilise le modèle \textbf{corotationnel} (\emph{fixed corotated}, Stomakhin et al. 2012), le
standard des solveurs MPM graphiques car il reste stable en grandes rotations et même en cas
d'inversion locale de $\mathbf F$ :
\[
  \mathbf P = 2\mu\,(\mathbf F - \mathbf R) + \lambda\,J\,(J-1)\,\mathbf F^{-\mathsf T}.
\]
$\mathbf P$ est la contrainte de Piola-Kirchhoff (par unité de surface de repos) ; celle qui
intervient dans la force nodale est $\bm\tau = \mathbf P\mathbf F^{\mathsf T}$ (Kirchhoff), qui se
simplifie en
\[
  \boxed{\;\bm\tau = 2\mu\,(\mathbf F - \mathbf R)\,\mathbf F^{\mathsf T} + \lambda\,J\,(J-1)\,\mathbf I\;}
\]
Le premier terme est la pénalité de \emph{distorsion} (écart entre $\mathbf F$ et sa rotation),
le second celle de \emph{volume} : c'est exactement la loi de pression du fluide,
$E(J-1)\mathbf I$, avec $\lambda J$ à la place de $E$. Les coefficients de Lamé viennent du module
d'Young $E$ et du coefficient de Poisson $\nu$ :
\[
  \mu = \frac{E}{2(1+\nu)},\qquad \lambda = \frac{E\nu}{(1+\nu)(1-2\nu)}.
\]
Une alternative classique est le Néo-Hookéen, $\bm\tau = \mu(\mathbf F\mathbf F^{\mathsf T} - \mathbf I) + \lambda\ln J\,\mathbf I$,
qui a les mêmes petites déformations mais explose si $J\to0$ ; on le mentionne pour l'exercice.

\subsection{La force nodale MLS-MPM}

Avec la contrainte de Cauchy $\bm\sigma = \bm\tau/J$ et le volume courant $V_p = J V_p^0$, la force
exercée par la particule sur le n\oe ud $i$ est
\[
  \mathbf f_i = -\sum_p V_p\,\bm\sigma_p\,\nabla w_{ip} = -\sum_p V_p^0\,\bm\tau_p\,\nabla w_{ip},
\]
et avec l'approximation MLS du gradient des poids, $\nabla w_{ip} \approx \frac{4}{\Delta x^2}\,w_{ip}\,(x_i - x_p)$,
\[
  \Delta t\,\mathbf f_i = w_{ip}\;\underbrace{\Big(-\Delta t\,\frac{4 V_p^0}{\Delta x^2}\,\bm\tau_p\Big)}_{\text{matrice } 2\times2}\;(x_i - x_p).
\]
Comparez au fluide : \code{stress * dpos} avec \code{stress} \emph{scalaire}, car
$\bm\tau = E(J-1)\mathbf I$ y est isotrope. Ici c'est une matrice appliquée à \code{dpos}.
Comme le terme APIC $m_p C_p\,(x_i - x_p)$ a exactement la même forme, on additionne les deux
matrices en une seule, appelée \code{affine} dans le code :
\[
  (mv)_i \mathrel{+}= w_{ip}\Big( m_p v_p + \underbrace{\big(-\Delta t\,\tfrac{4V_p^0}{\Delta x^2}\bm\tau_p + m_p C_p\big)}_{\texttt{affine}}\,(x_i - x_p)\Big).
\]

\subsection{L'endommagement}
\label{sec:damage}

\paragraph{Critère.} Un pont casse en \emph{traction} (fibre inférieure à mi-portée, fibre supérieure
aux appuis), pas en compression. On mesure donc l'\textbf{allongement principal maximal}
\[
  \varepsilon_p = \max(s_1, s_2) - 1,
\]
directement lu sur la SVD de $\mathbf F_p$. C'est un critère de type Rankine (contrainte principale
maximale) écrit en déformation. Le critère de Von Mises,
$\sigma_{\rm eq} = \sqrt{\tfrac32\,\mathbf s:\mathbf s}$ avec $\mathbf s$ le déviateur de $\bm\sigma$,
est le bon choix pour la \emph{plasticité} des métaux, mais il est symétrique traction/compression :
un pilier comprimé s'endommagerait autant qu'un tablier tendu. Pour un matériau fragile
(béton, pierre, verre) le critère en traction est plus fidèle.

\paragraph{Loi d'évolution.} Deux seuils, $\varepsilon_0$ (début de l'endommagement) et
$\varepsilon_f$ (rupture), et une interpolation linéaire entre les deux :
\[
  D_p \leftarrow \max\!\Big(D_p,\ \operatorname{clamp}\Big(\frac{\varepsilon_p - \varepsilon_0}{\varepsilon_f - \varepsilon_0},\,0,\,1\Big)\Big).
\]
Le $\max$ rend $D_p$ \textbf{irréversible} : une fissure ne se referme pas. La raideur effective
est
\[
  E_{\rm eff} = (1 - D_p)\,E \quad\Longleftrightarrow\quad \mu_{\rm eff} = (1-D_p)\mu,\ \lambda_{\rm eff} = (1-D_p)\lambda.
\]
C'est un \emph{adoucissement} : plus la particule est endommagée, moins elle résiste, donc plus elle
s'allonge, donc plus $D_p$ croît. La déformation se localise d'elle-même dans une bande étroite :
c'est la fissure.

\subsection{La rupture : supprimer la cohésion}
\label{sec:rupture}

Marquer simplement \code{broken = 1} et mettre la contrainte à zéro ne suffit pas : la particule
continuerait de partager des n\oe uds avec ses voisines et resterait \og collée \fg{} au reste
par le champ de vitesse commun. Il faut lui donner un comportement de matériau \textbf{sans
cohésion} (du sable) : il résiste à la compression, jamais à la traction. Concrètement, à chaque
pas et pour chaque particule rompue :
\begin{enumerate}
  \item à l'instant de la rupture, $\mathbf F_p \leftarrow \mathbf I$ : la particule oublie la
        déformation accumulée (l'énergie élastique de traction est perdue, comme dans une vraie
        fissure) ;
  \item après la mise à jour $\mathbf F_p \leftarrow (\mathbf I + \Delta t C_p)\mathbf F_p$, on écrête
        les allongements principaux :
        \[
          s_k \leftarrow \operatorname{clamp}(s_k,\ 0.1,\ 1),\qquad \mathbf F_p \leftarrow \mathbf U\,\bm\Sigma\,\mathbf V^{\mathsf T}.
        \]
        Un étirement $s_k > 1$ est immédiatement ramené à 1 : aucune traction ne peut naître.
        La borne inférieure évite l'inversion sous un écrasement extrême ;
  \item la contrainte est calculée avec la raideur \emph{pleine} ($k = 1$, pas $1-D$), sinon la
        particule ne résisterait même plus à la compression et les morceaux s'interpénétreraient.
\end{enumerate}
C'est le même mécanisme que la plasticité de la neige de Stomakhin et al. (2013), poussé à
l'extrême (aucune traction admise). Il ne coûte qu'une SVD par particule rompue.

\paragraph{Limite à connaître.} Deux morceaux séparés par moins d'une cellule interagissent
encore à travers les n\oe uds partagés (ils \og se voient \fg{} par la grille). Les fissures ont
donc une largeur minimale d'environ une cellule, et deux fragments en contact glissent sans
frottement bien défini. Les remèdes classiques sont CPIC (\emph{compatible PIC} : deux jeux de
vitesses nodales de part et d'autre d'une fissure) ou un réseau explicite de liaisons entre
particules initialement voisines (cassées quand leur allongement
$\epsilon_{pq} = |\mathbf x_p - \mathbf x_q|/d^0_{pq} - 1$ dépasse un seuil). Pour un premier
solveur, le modèle continu ci-dessus donne déjà des fissures nettes.

\subsection{Stabilité et choix des paramètres}
\label{sec:stab}

Le schéma est explicite : $\Delta t < \Delta x / c_p$ avec la vitesse des ondes de compression
\[
  c_p = \sqrt{\frac{\lambda + 2\mu}{\rho_s}}.
\]
Dans le code, $\Delta t = 0.4\,\Delta x/c_p$ est \emph{calculé} à partir des paramètres, pas
choisi à la main. Valeurs de la partie I :

\begin{center}
\begin{tabular}{@{}lll@{}}
\toprule
Paramètre & Valeur & Rôle \\
\midrule
$n$ & 128 & n\oe uds par côté, $\Delta x = 1/128$ \\
$\rho_s$ & 2 & masse volumique (2 fois l'eau : les morceaux coulent) \\
$E$, $\nu$ & 3000, 0.3 & $\mu = 1154$, $\lambda = 1731$, $c_p \approx 45$ \\
$\Delta t$ & $0.4\,\Delta x/c_p \approx 7\cdot10^{-5}$ & calculé \\
$\varepsilon_0$, $\varepsilon_f$ & 0.05, 0.20 & seuils d'endommagement et de rupture \\
$V_p$ & $(\Delta x/2)^2$ & 4 particules par cellule, sur un réseau régulier \\
poutre & $[0.15, 0.85]\times[0.55, 0.61]$ & extrémités noyées dans des piliers de largeur 0.06 \\
\bottomrule
\end{tabular}
\end{center}

$E = 3000$ est \emph{très} mou pour un vrai matériau ($\sim 30$~GPa pour du béton), mais c'est ce
qui rend la flèche visible à l'\oe il ; de même $\varepsilon_f = 0.2$ est ductile. L'ordre de
grandeur qui compte est la déformation de flexion d'une poutre encastrée sous son poids,
$\varepsilon \approx \rho_s g L^2 / (2 h E)$ : avec $L = 0.58$, $h = 0.06$, on trouve
$\varepsilon \approx 0.02$ par unité de charge, à comparer aux seuils. Le curseur de charge
multiplie $g$ et permet de balayer $\varepsilon$ de 0 à $0.1$.

% ================================================================
\part{Le solide seul}

% ================================================================
\section{Étape 1 : champs et loi de comportement}

Une particule solide porte, en plus de $x, v, C$ : $\mathbf F$ (matrice $2\times2$), $D$ (réel)
et \code{broken} (entier). La grille est la même que pour le fluide, plus un masque d'obstacles
qui servira à encastrer la poutre.

\begin{lstlisting}
x_s = ti.Vector.field(2, ti.f32, n_solid)      # position
v_s = ti.Vector.field(2, ti.f32, n_solid)      # vitesse
C_s = ti.Matrix.field(2, 2, ti.f32, n_solid)   # matrice affine APIC (~ grad v)
F_s = ti.Matrix.field(2, 2, ti.f32, n_solid)   # gradient de déformation (I au repos)
D_s = ti.field(ti.f32, n_solid)                # endommagement dans [0, 1]
broken_s = ti.field(ti.i32, n_solid)           # 1 si rompue
col_s = ti.Vector.field(3, ti.f32, n_solid)    # couleur d'affichage

grid_v = ti.Vector.field(2, ti.f32, (n_grid, n_grid))
grid_m = ti.field(ti.f32, (n_grid, n_grid))
mask = ti.field(ti.i32, (n_grid, n_grid))      # 1 = pilier rigide
\end{lstlisting}

La loi de comportement est une fonction Taichi (\code{@ti.func}) appelable depuis un kernel.
Elle reçoit $\mathbf F$ et les coefficients de Lamé (éventuellement déjà réduits par $1-D$) et
renvoie $\bm\tau$ :

\begin{lstlisting}
@ti.func
def kirchhoff_stress(F, mu: float, la: float):
    U, sig, V = ti.svd(F)
    R = U @ V.transpose()                       # rotation de la décomposition polaire
    J = F.determinant()
    I = ti.Matrix.identity(ti.f32, 2)
    return 2.0 * mu * (F - R) @ F.transpose() + la * J * (J - 1.0) * I
\end{lstlisting}

C'est la formule encadrée du \S\ref{sec:loi}, ligne pour ligne. Pour essayer le Néo-Hookéen,
remplacez le \code{return} par \code{mu * (F @ F.transpose() - I) + la * ti.log(J) * I}.

% ================================================================
\section{Étape 2 : P2G solide}

Même squelette que le \code{P2G} du fluide : stencil, poids, puis dépôt sur les 9 n\oe uds. Deux
différences : la contrainte est une matrice, et le kernel \emph{ne remet pas la grille à zéro}
(pour pouvoir s'ajouter au dépôt du fluide en partie II ; en partie I, un petit kernel
\code{clear_grid} s'en charge avant).

\begin{lstlisting}
@ti.kernel
def P2G_solid(grid_m: ti.template(), grid_v: ti.template(),
              x: ti.template(), v: ti.template(), C: ti.template(), F: ti.template(),
              D: ti.template(), broken: ti.template(),
              inv_dx: float, dt: float, dx: float,
              mu: float, la: float, p_mass: float, p_vol: float):
    for p in x:
        base = (x[p] * inv_dx - 0.5).cast(int)
        fx = x[p] * inv_dx - base
        w = [0.5 * (1.5 - fx)**2, 0.75 - (fx - 1.0)**2, 0.5 * (fx - 0.5)**2]

        k = 1.0 - D[p]                 # raideur effective (1 - D) E
        if broken[p] == 1:
            k = 1.0                    # rompue : raideur pleine, F écrêté dans G2P
        tau = kirchhoff_stress(F[p], k * mu, k * la)

        affine = (-dt * p_vol * 4.0 * inv_dx * inv_dx) * tau + p_mass * C[p]

        for i, j in ti.static(ti.ndrange(3, 3)):
            offset = ti.Vector([i, j])
            d_pos = (offset - fx) * dx
            weight = w[i].x * w[j].y
            grid_v[base + offset] += weight * (p_mass * v[p] + affine @ d_pos)
            grid_m[base + offset] += weight * p_mass
\end{lstlisting}

\paragraph{Ligne par ligne.}
\begin{itemize}
  \item \code{k} : le facteur $(1-D_p)$ du \S\ref{sec:damage}, forcé à 1 pour une particule
        rompue (\S\ref{sec:rupture}, point 3).
  \item \code{affine} : la matrice $-\Delta t\,\frac{4V_p^0}{\Delta x^2}\bm\tau_p + m_p C_p$. Notez
        que \code{p_vol} est le volume de \emph{repos} $V_p^0$ : c'est $\bm\tau$ (et non $\bm\sigma$)
        qui va avec.
  \item \code{affine @ d_pos} remplace le \code{C[p] @ d_pos} et le \code{stress * d_pos} du fluide
        d'un seul coup.
\end{itemize}

% ================================================================
\section{Étape 3 : grille et encastrement}

La mise à jour de la grille est \emph{identique} à celle du fluide : on la recopie dans
\code{mpm_solid.py} sous le nom \code{grid_update} pour que la partie I soit autonome (en partie II
on utilisera directement \code{grid_step} du fluide).

\begin{lstlisting}
@ti.kernel
def grid_update(grid_m: ti.template(), grid_v: ti.template(), mask: ti.template(),
                dt: float, g: float, bound: int, n_grid: int):
    for i, j in grid_m:
        if grid_m[i, j] > 0:
            grid_v[i, j] /= grid_m[i, j]
            grid_v[i, j].y -= dt * g
            # parois glissantes du domaine (comme le fluide) ...
            if mask[i, j] == 1:
                grid_v[i, j] = [0.0, 0.0]
\end{lstlisting}

\paragraph{Encastrement.} Un n\oe ud du masque a une vitesse nulle. Une particule dont les 9
n\oe uds du stencil sont dans le masque ne bouge donc jamais : c'est un encastrement parfait, sans
aucun code supplémentaire. Il suffit de \emph{noyer} les extrémités de la poutre dans les
piliers :

\begin{lstlisting}
@ti.kernel
def init_mask():
    for i, j in mask:
        pos = ti.Vector([i, j]) * dx
        left = beam_x0 <= pos.x <= beam_x0 + pillar_w
        right = beam_x1 - pillar_w <= pos.x <= beam_x1
        mask[i, j] = 1 if ((left or right) and pos.y <= beam_y1) else 0
\end{lstlisting}

Les piliers montent jusqu'au \emph{dessus} de la poutre (\code{pos.y <= beam_y1}) : les particules
des extrémités sont entièrement gelées. Pour un appui simple (rotation libre), on arrêterait les
piliers sous la poutre.

% ================================================================
\section{Étape 4 : G2P solide, version élastique}

La première moitié est le \code{G2P} du fluide (vitesse et $C_p$). Ensuite, à la place de
$J_p \leftarrow J_p(1+\Delta t\,\tr C_p)$, on met à jour $\mathbf F_p$. Voici le kernel complet ;
les deux blocs \code{if broken} / \code{elif use_damage} sont expliqués aux étapes suivantes
et sont inactifs en mode élastique.

\begin{lstlisting}
@ti.kernel
def G2P_solid(grid_v: ti.template(),
              x: ti.template(), v: ti.template(), C: ti.template(), F: ti.template(),
              D: ti.template(), broken: ti.template(),
              inv_dx: float, dt: float, dx: float, bound: int,
              eps0: float, epsf: float, use_damage: int, use_rupture: int):
    I = ti.Matrix.identity(ti.f32, 2)
    for p in x:
        base = (x[p] * inv_dx - 0.5).cast(int)
        fx = x[p] * inv_dx - base
        w = [0.5 * (1.5 - fx)**2, 0.75 - (fx - 1.0)**2, 0.5 * (fx - 0.5)**2]
        v_new = ti.Vector.zero(ti.f32, 2)
        C_new = ti.Matrix.zero(ti.f32, 2, 2)
        for i, j in ti.static(ti.ndrange(3, 3)):
            offset = ti.Vector([i, j])
            d_pos = (offset - fx) * dx
            weight = w[i].x * w[j].y
            g_v = grid_v[base + offset]
            v_new += weight * g_v
            C_new += weight * g_v.outer_product(d_pos) * (4.0 * inv_dx * inv_dx)
        v[p] = v_new
        C[p] = C_new

        # --- gradient de déformation : dF/dt = (grad v) F, grad v ~ C
        F_new = (I + dt * C_new) @ F[p]

        if broken[p] == 1:                          # étape 6 : rupture
            U, sig, V = ti.svd(F_new)
            for d in ti.static(range(2)):
                sig[d, d] = ti.math.clamp(sig[d, d], 0.1, 1.0)
            F_new = U @ sig @ V.transpose()

        elif use_damage == 1:                       # étape 5 : endommagement
            U, sig, V = ti.svd(F_new)
            eps = ti.max(sig[0, 0], sig[1, 1]) - 1.0
            D_new = ti.math.clamp((eps - eps0) / (epsf - eps0), 0.0, 1.0)
            D[p] = ti.max(D[p], D_new)
            if use_rupture == 1 and D[p] >= 1.0:
                broken[p] = 1
                F_new = I

        F[p] = F_new
        x[p] += dt * v[p]
        x[p] = ti.math.clamp(x[p], bound * dx, 1.0 - bound * dx)
\end{lstlisting}

\paragraph{Ce que l'on doit voir (touche 1, charge 1).} La poutre fléchit d'environ une cellule
à mi-portée, oscille quelques fois et se stabilise (la dissipation numérique de PIC/APIC amortit
les vibrations). Avec le curseur de charge à 3, la flèche triple, sans jamais casser : un solide
élastique pur revient toujours. En mode couleur \og déformation \fg{} (Espace), la fibre
inférieure à mi-portée et les fibres supérieures aux encastrements apparaissent en rouge
(traction), les fibres opposées en bleu (compression) : c'est le diagramme des moments d'une
poutre bi-encastrée, dessiné par la simulation.

% ================================================================
\section{Étape 5 : endommagement}

Le bloc \code{elif use_damage == 1} est la loi du \S\ref{sec:damage}, mot pour mot :
\begin{itemize}
  \item \code{eps} : allongement principal maximal, lu sur la SVD du \emph{nouveau} $\mathbf F$ ;
  \item \code{D_new} : interpolation linéaire entre $\varepsilon_0$ et $\varepsilon_f$, bornée à $[0,1]$ ;
  \item \code{ti.max(D[p], D_new)} : irréversibilité.
\end{itemize}
Et $D_p$ est consommé dans \code{P2G_solid} via \code{k = 1 - D[p]}. C'est tout : l'endommagement
tient en quatre lignes.

\paragraph{Ce que l'on doit voir (touche 2).} À charge 1, rien ne change ($\varepsilon < \varepsilon_0$,
$D_{\max} = 0$ dans le panneau). À charge 2, les particules jaunissent aux encastrements (dessus)
et à mi-portée (dessous), $D_{\max}$ monte vers $0.3$ et la flèche augmente un peu : la poutre est
fissurée mais tient. Ramenez la charge à 1 : la flèche diminue mais la couleur reste, $D$ ne
redescend jamais. À charge 3, $D_{\max}$ atteint 1 par endroits et la poutre s'affaisse sans se
séparer : les particules de raideur nulle restent \og accrochées \fg{} par la grille. C'est le
problème annoncé au \S\ref{sec:rupture}, et la raison de l'étape suivante.

% ================================================================
\section{Étape 6 : rupture}

Deux morceaux de code. D'abord la \emph{transition}, dans le bloc d'endommagement :
\begin{lstlisting}
            if use_rupture == 1 and D[p] >= 1.0:
                broken[p] = 1
                F_new = I
\end{lstlisting}
Ensuite le comportement \emph{permanent} d'une particule rompue, en tête du \code{if} :
\begin{lstlisting}
        if broken[p] == 1:
            U, sig, V = ti.svd(F_new)
            for d in ti.static(range(2)):
                sig[d, d] = ti.math.clamp(sig[d, d], 0.1, 1.0)
            F_new = U @ sig @ V.transpose()
\end{lstlisting}
La boucle \code{for d in ti.static(range(2))} est déroulée à la compilation et écrête les deux
allongements principaux ; \code{U @ sig @ V.transpose()} recompose $\mathbf F$ avec la même
rotation. Une particule rompue ne repasse jamais par le bloc d'endommagement (\code{elif}).

\paragraph{Ce que l'on doit voir (touche 3).} Chargez progressivement. Vers 2.5 à 3, les premières
particules rouges apparaissent \emph{aux encastrements}, sur la fibre supérieure : c'est là que le
moment fléchissant d'une poutre bi-encastrée est maximal ($qL^2/12$ contre $qL^2/24$ à
mi-portée). La poutre se détache des piliers, devient une poutre sur appuis, puis se fissure
\emph{à mi-portée} par le dessous : trois rotules, mécanisme de ruine classique. Les deux moitiés
tombent, tournent, heurtent le sol et s'y brisent en éclats (l'impact à plusieurs mètres par
seconde sur une paroi rigide impose des déformations bien supérieures à $\varepsilon_f$).

Mesures faites sur ce code, avec les paramètres de la table du \S\ref{sec:stab}, sur 1~s :

\begin{center}
\begin{tabular}{@{}llll@{}}
\toprule
Charge & Modèle & $D_{\max}$ & Comportement \\
\midrule
1 & 3 & 0 & flèche $\approx 0.014$, aucune particule endommagée \\
2 & 3 & 0.28 & fissuration visible, la poutre tient \\
3 & 3 & 1 & rupture aux encastrements à $t\approx0.1$~s, à mi-portée vers $t\approx0.2$~s, chute \\
3 & 1 & -- & (élastique) flèche $\approx 0.04$, oscille, ne casse pas \\
\bottomrule
\end{tabular}
\end{center}

% ================================================================
\section{Étape 7 : le programme de démonstration}

\code{MpM/demo_solid.py} assemble le tout. Les particules sont placées sur un réseau régulier de
$2\times2$ par cellule (plus régulier qu'un tirage aléatoire, ce qui compte pour un solide dont on
regarde la forme) :

\begin{lstlisting}
@ti.kernel
def init_beam(x, v, C, F, D, broken, x0: float, y0: float, n_px: int, spacing: float):
    for p in x:
        i = p % n_px
        j = p // n_px
        x[p] = [x0 + (i + 0.5) * spacing, y0 + (j + 0.5) * spacing]
        v[p] = [0.0, 0.0]
        C[p] = ti.Matrix.zero(ti.f32, 2, 2)
        F[p] = ti.Matrix.identity(ti.f32, 2)
        D[p] = 0.0
        broken[p] = 0
\end{lstlisting}
(les arguments \code{x, v, ...} sont des \code{ti.template()} dans le fichier ; on abrège ici.)

Un pas de temps, avec la charge \code{load} comme multiplicateur de gravité et les deux drapeaux :
\begin{lstlisting}
def substep(load, use_damage, use_rupture):
    clear_grid()
    P2G_solid(grid_m, grid_v, x_s, v_s, C_s, F_s, D_s, broken_s,
              inv_dx, dt, dx, mu_s, la_s, p_mass, p_vol)
    grid_update(grid_m, grid_v, mask, dt, load * g, bound, n_grid)
    G2P_solid(grid_v, x_s, v_s, C_s, F_s, D_s, broken_s,
              inv_dx, dt, dx, bound, eps0, epsf, use_damage, use_rupture)
\end{lstlisting}

\paragraph{Affichage.} On ne peut plus se contenter d'une image de la grille : il faut voir les
particules solides et leur couleur. GGUI dessine des disques depuis un champ de positions avec
\code{canvas.circles(x_s, radius=..., per_vertex_color=col_s)}, par-dessus l'image de fond
(piliers en gris). Le kernel \code{solid_colors} remplit \code{col_s} : gris $\to$ jaune selon $D$,
rouge si rompue ; ou bleu $\to$ rouge selon l'allongement principal (mode déformation).
Le curseur de charge est un \code{gui.slider_float}. Le code complet est en annexe~\ref{app:demo}.

\paragraph{Expériences.}
\begin{itemize}
  \item Matériau fragile : $\varepsilon_0 = 0.02$, $\varepsilon_f = 0.03$. La rupture devient
        brutale, sans phase jaune visible.
  \item Matériau ductile : $\varepsilon_f = 0.5$. La poutre pend comme du caoutchouc avant de céder.
  \item Appui simple : piliers arrêtés sous la poutre (\code{pos.y <= beam_y0}). La rupture se
        produit alors d'abord à mi-portée.
  \item Console : supprimez le pilier droit. Une poutre encastrée-libre casse à l'encastrement.
  \item Critère de Von Mises : calculez $\bm\sigma = \bm\tau/J$ dans \code{G2P_solid}, puis
        $\sigma_{\rm eq}$, et remplacez \code{eps} par $\sigma_{\rm eq}/E$. Observez que les zones
        comprimées s'endommagent aussi.
\end{itemize}

% ================================================================
\part{Couplage avec le fluide APIC}

% ================================================================
\section{Le principe : une grille, deux phases}

\subsection{Pourquoi c'est presque gratuit}

Le fluide dépose $m_i$, $(mv)_i$ et sa force de pression ; le solide dépose $m_i$, $(mv)_i$ et sa
force élastique. Sur un n\oe ud partagé,
\[
  m_i = m_i^{f} + m_i^{s},\qquad (mv)_i = (mv)_i^{f} + (mv)_i^{s} + \Delta t\,(\mathbf f_i^{f} + \mathbf f_i^{s}),
  \qquad v_i = (mv)_i / m_i,
\]
et chaque phase relit la \emph{même} vitesse $v_i$. C'est un \textbf{couplage monolithique} : une
seule équation de quantité de mouvement pour le mélange, résolue une fois par pas de temps. Il
n'y a ni force de contact à calculer, ni itération entre les deux solveurs. Les conséquences
physiques :
\begin{itemize}
  \item \textbf{Pas d'interpénétration} : là où l'eau pousse le tablier, les n\oe uds portent les deux
        masses et une seule vitesse ; la pression du fluide accélère le solide et la raideur du
        solide freine le fluide. C'est exactement la réaction hydrodynamique voulue.
  \item \textbf{Adhérence à l'interface} : une vitesse commune signifie condition de \emph{no-slip}
        entre eau et solide. Pour un tablier de pont c'est parfaitement acceptable.
  \item \textbf{Mélange sur une cellule} : l'interface n'est pas nette, elle a l'épaisseur du
        stencil. Un fluide et un solide qui se croisent à grande vitesse échangent un peu trop de
        quantité de mouvement. C'est la limite connue de cette approche ; le remède (deux champs
        de vitesse nodaux et une condition de contact entre eux) fait l'objet du dernier
        paragraphe.
\end{itemize}

\subsection{L'ordre des kernels}

\begin{center}
\begin{tikzpicture}[
  stp/.style={draw, rounded corners, minimum width=3.9cm, minimum height=0.8cm, align=center, font=\small},
  arr/.style={-{Latex[length=2.5mm]}, thick}]
  \node[stp,fill=blue!5] (a) {1. \code{P2G} fluide\\ \emph{remet la grille à zéro}};
  \node[stp,fill=orange!10, right=0.5cm of a] (b) {2. \code{P2G_solid}\\ s'ajoute};
  \node[stp,fill=gray!10, right=0.5cm of b] (c) {3. \code{grid_step}\\ $v_i$, gravité, parois, piliers};
  \node[stp,fill=blue!5, below=0.6cm of c] (d) {4. \code{G2P} fluide\\ $v, C, J$};
  \node[stp,fill=orange!10, left=0.5cm of d] (e) {5. \code{G2P_solid}\\ $v, C, \mathbf F, D$, advection};
  \node[stp,fill=blue!5, left=0.5cm of e] (f) {6. \code{advect_fluid}\\ $x \mathrel{+}= \Delta t\, v$};
  \draw[arr] (a) -- (b); \draw[arr] (b) -- (c); \draw[arr] (c) -- (d);
  \draw[arr] (d) -- (e); \draw[arr] (e) -- (f);
\end{tikzpicture}
\end{center}

\begin{itemize}
  \item Le \code{P2G} du fluide remet la grille à zéro avant de déposer : il doit donc passer
        \emph{en premier}, et \code{P2G_solid} (qui ne remet rien à zéro) s'ajoute ensuite. Aucune
        modification de \code{APIC/APIC.py} n'est nécessaire.
  \item \code{grid_step} du fluide sert tel quel : son masque \code{solid} (obstacles) devient le
        masque des piliers.
  \item Le \code{G2P} du fluide n'avance pas les particules (c'était le rôle de
        \code{time_integration}) ; un kernel \code{advect_fluid} de trois lignes le fait avec le
        \emph{même} $\Delta t$.
\end{itemize}

\subsection{Un seul pas de temps}

Le fluide impose $\Delta t < \Delta x / c_f$ avec $c_f = \sqrt{E_f/\rho_f} = 20$, le solide
$\Delta t < \Delta x / c_p$ avec $c_p \approx 45$. On prend
\[
  \Delta t = 0.4\,\frac{\Delta x}{\max(c_f, c_p)} \approx 7\cdot10^{-5}\quad (n = 128),
\]
et on l'utilise \emph{partout} : dans les deux P2G (facteur $\Delta t$ devant les forces), dans
les deux G2P (mise à jour de $J$ et de $\mathbf F$) et dans les deux advections.

\paragraph{Remarque sur \code{main.py}.} Dans le programme fluide actuel, le \code{dt} fixe passé
à \code{P2G} et \code{G2P} n'est pas celui, adaptatif, que \code{time_integration} utilise pour
avancer les particules. Les forces et l'advection ne sont alors pas intégrées avec le même pas,
ce qui fausse $J_p$ et la pression. Le solveur couplé corrige ce point en calculant un $\Delta t$
unique à partir de la condition CFL \emph{avant} la simulation. Si l'on tient à un pas adaptatif,
il faut le calculer en début de sous-pas et le passer à tous les kernels.

% ================================================================
\section{Le programme couplé : \texttt{main\_fsi.py}}

\subsection{Scène}

Deux piliers (masque), un tablier encastré entre eux à mi-hauteur, un bassin au fond, et un bloc
d'eau lâché au-dessus du tablier :

\begin{center}
\begin{tikzpicture}[scale=5]
  \draw[thick] (0,0) rectangle (1,1);
  \fill[blue!30] (0.03,0.03) rectangle (0.97,0.22);
  \fill[blue!30] (0.30,0.62) rectangle (0.70,0.95);
  \fill[gray!60] (0.15,0) rectangle (0.21,0.46);
  \fill[gray!60] (0.79,0) rectangle (0.85,0.46);
  \fill[orange!60] (0.15,0.40) rectangle (0.85,0.46);
  \node[font=\small] at (0.5,0.785) {bloc d'eau};
  \node[font=\small] at (0.5,0.125) {bassin};
  \node[font=\small] at (0.5,0.43) {tablier};
  \node[font=\small, rotate=90] at (0.18,0.23) {pilier};
  \node[font=\small, rotate=90] at (0.82,0.23) {pilier};
  \draw[-{Latex}, thick] (0.5,0.60) -- (0.5,0.50);
\end{tikzpicture}
\end{center}

\begin{lstlisting}
import taichi as ti
from APIC.APIC import P2G, G2P, grid_step
from MpM.mpm_solid import (P2G_solid, G2P_solid, init_beam, solid_colors, solid_stats)

ti.init(arch=ti.gpu)

n_grid = 128
dx = 1.0 / n_grid
inv_dx = float(n_grid)
bound = 3
render_substep = 20

# fluide (comme main.py)
rho_f, E_f = 1.0, 400.0
p_vol_f = (0.5 * dx) ** 2
p_mass_f = p_vol_f * rho_f
water_x0, water_x1, water_y0, water_y1 = 0.30, 0.70, 0.62, 0.95   # bloc qui tombe
pool_x0, pool_x1, pool_y0, pool_y1 = 0.03, 0.97, 0.03, 0.22        # bassin
n_water = int(4 * (water_x1 - water_x0) * (water_y1 - water_y0) * n_grid * n_grid)
n_pool = int(4 * (pool_x1 - pool_x0) * (pool_y1 - pool_y0) * n_grid * n_grid)
n_fluid = n_water + n_pool

# solide (comme demo_solid.py)
rho_s, E_s, nu_s = 2.0, 3000.0, 0.3
mu_s = E_s / (2 * (1 + nu_s))
la_s = E_s * nu_s / ((1 + nu_s) * (1 - 2 * nu_s))
eps0, epsf = 0.05, 0.20
use_damage, use_rupture = 1, 1
p_spacing = 0.5 * dx
p_vol_s = p_spacing ** 2
p_mass_s = p_vol_s * rho_s
beam_x0, beam_x1, beam_y0, beam_y1 = 0.15, 0.85, 0.40, 0.46
pillar_w = 0.06
n_px = int(round((beam_x1 - beam_x0) / p_spacing))
n_py = int(round((beam_y1 - beam_y0) / p_spacing))
n_solid = n_px * n_py

# pas de temps commun
g = 9.81
c_f = (E_f / rho_f) ** 0.5
c_p = ((la_s + 2 * mu_s) / rho_s) ** 0.5
dt = 0.4 * dx / max(c_f, c_p)
\end{lstlisting}

\subsection{Champs}

Les champs du fluide portent le suffixe \code{_f}, ceux du solide \code{_s} ; la grille est
unique. Le masque s'appelle \code{solid} parce que c'est le nom attendu par \code{grid_step}
(et il est en \code{ti.f32} comme dans \code{main.py} ; la comparaison \code{== 1} fonctionne).

\begin{lstlisting}
x_f = ti.Vector.field(2, ti.f32, n_fluid)
v_f = ti.Vector.field(2, ti.f32, n_fluid)
C_f = ti.Matrix.field(2, 2, ti.f32, n_fluid)
J_f = ti.field(ti.f32, n_fluid)

x_s = ti.Vector.field(2, ti.f32, n_solid)
v_s = ti.Vector.field(2, ti.f32, n_solid)
C_s = ti.Matrix.field(2, 2, ti.f32, n_solid)
F_s = ti.Matrix.field(2, 2, ti.f32, n_solid)
D_s = ti.field(ti.f32, n_solid)
broken_s = ti.field(ti.i32, n_solid)
col_s = ti.Vector.field(3, ti.f32, n_solid)

grid_v = ti.Vector.field(2, ti.f32, (n_grid, n_grid))
grid_m = ti.field(ti.f32, (n_grid, n_grid))
solid = ti.field(ti.f32, (n_grid, n_grid))     # 1 = pilier
img = ti.Vector.field(3, ti.f32, (n_grid, n_grid))
\end{lstlisting}

\subsection{Initialisation et advection du fluide}

\begin{lstlisting}
@ti.kernel
def init_fluid():
    for p in x_f:
        if p < n_water:
            x_f[p] = [water_x0 + ti.random() * (water_x1 - water_x0),
                      water_y0 + ti.random() * (water_y1 - water_y0)]
        else:
            x_f[p] = [pool_x0 + ti.random() * (pool_x1 - pool_x0),
                      pool_y0 + ti.random() * (pool_y1 - pool_y0)]
        v_f[p] = [0.0, 0.0]
        C_f[p] = ti.Matrix.zero(ti.f32, 2, 2)
        J_f[p] = 1.0
    for i, j in solid:
        pos = ti.Vector([i, j]) * dx
        left = beam_x0 <= pos.x <= beam_x0 + pillar_w
        right = beam_x1 - pillar_w <= pos.x <= beam_x1
        solid[i, j] = 1.0 if ((left or right) and pos.y <= beam_y1) else 0.0

@ti.kernel
def advect_fluid():
    for p in x_f:
        x_f[p] += dt * v_f[p]
        x_f[p] = ti.math.clamp(x_f[p], bound * dx, 1.0 - bound * dx)

def init_scene():
    init_fluid()
    init_beam(x_s, v_s, C_s, F_s, D_s, broken_s, beam_x0, beam_y0, n_px, p_spacing)
\end{lstlisting}

\subsection{Le pas de temps couplé}

C'est le c\oe ur de la partie II, et il tient en six appels dans l'ordre du schéma :

\begin{lstlisting}
def substep():
    P2G(grid_m, grid_v, x_f, v_f, C_f, J_f, inv_dx, dt, dx, E_f, p_mass_f, p_vol_f)   # 1
    P2G_solid(grid_m, grid_v, x_s, v_s, C_s, F_s, D_s, broken_s,
              inv_dx, dt, dx, mu_s, la_s, p_mass_s, p_vol_s)                           # 2
    grid_step(grid_m, grid_v, solid, dt, g, bound, n_grid)                             # 3
    G2P(grid_m, grid_v, x_f, v_f, C_f, J_f, inv_dx, dt, dx, E_f, p_mass_f, p_vol_f)   # 4
    G2P_solid(grid_v, x_s, v_s, C_s, F_s, D_s, broken_s,
              inv_dx, dt, dx, bound, eps0, epsf, use_damage, use_rupture)              # 5
    advect_fluid()                                                                     # 6
\end{lstlisting}

\paragraph{Ce qui se passe sur un n\oe ud sous le tablier quand l'eau arrive.} Après l'appel 1, le
n\oe ud porte la masse et la quantité de mouvement de l'eau qui descend, plus la force de pression
(l'eau se comprime à l'impact, $J_f < 1$, \code{stress} $> 0$). Après l'appel 2, il porte en plus
la masse du tablier (au repos) et sa force élastique. L'appel 3 fait la moyenne pondérée par les
masses : le n\oe ud descend, moins vite que l'eau. À l'appel 4 l'eau relit cette vitesse : elle
est freinée (et $J_f$ diminue, la pression monte). À l'appel 5 le tablier relit la même vitesse :
il est poussé vers le bas, son $\mathbf F$ se met à fléchir, et au pas suivant sa force élastique
repoussera le n\oe ud vers le haut. L'équilibre dynamique s'établit en quelques pas.

\subsection{Affichage et boucle}

Fond gris/bleu nuit depuis le masque, puis deux appels à \code{canvas.circles} : l'eau en bleu
uni, le solide coloré par $D$ ou par la déformation :

\begin{lstlisting}
        stats = solid_stats(D_s, broken_s)
        solid_colors(F_s, D_s, broken_s, col_s, color_mode, epsf)
        canvas.set_image(img)
        canvas.circles(x_f, radius=0.5 * p_spacing, color=(0.35, 0.65, 1.0))
        canvas.circles(x_s, radius=0.6 * p_spacing, per_vertex_color=col_s)
\end{lstlisting}
Le reste de la boucle (événements, panneau de texte) est celui de \code{main.py} ; voir
l'annexe~\ref{app:fsi}.

\subsection{Ce que l'on doit voir}

Le bloc d'eau tombe pendant $0.2$~s et frappe le tablier à environ $2$~m/s. Le tablier fléchit
(en mode déformation : rouge dessous à mi-portée, rouge dessus aux piliers), jaunit, puis vers
$t \approx 0.3$~s les premières particules rouges apparaissent aux encastrements. Le tablier se
décroche des deux piliers, tombe d'un bloc dans le bassin en entraînant l'eau qui le recouvrait,
et coule (il est deux fois plus dense que l'eau). L'eau du bloc déferle sur les piliers et
rejoint le bassin. Mesures sur ce code ($n = 128$, environ $20\,000$ particules de fluide et
$2\,700$ de solide, 15~s de calcul par seconde simulée sur GPU) :

\begin{center}
\begin{tabular}{@{}lllll@{}}
\toprule
$t$ (s) & $D_{\max}$ & rompues & $y$ mi-portée & commentaire \\
\midrule
0.10 & 0 & 0 & 0.418 & chute libre de l'eau, tablier au repos sous son poids \\
0.20 & 0.20 & 0 & 0.398 & impact, flexion, début de l'endommagement \\
0.30 & 1 & 65 & 0.256 & rupture aux deux encastrements, le tablier tombe \\
0.50 & 1 & 272 & 0.299 & tablier dans le bassin, éclats aux extrémités \\
1.00 & 1 & 290 & 0.254 & tablier au fond, l'eau se calme \\
\bottomrule
\end{tabular}
\end{center}

La masse totale sur la grille (fluide $+$ solide) est conservée à la précision machine, $J_f$
reste dans $[0.65, 1.25]$ (les valeurs extrêmes sont des particules d'eau coincées entre le
tablier qui tombe et un pilier) et la simulation reste stable indéfiniment. À $n = 256$ le
programme tourne tel quel : $\Delta t$ est divisé par deux automatiquement.

\paragraph{Expériences.}
\begin{itemize}
  \item Rendez le tablier plus résistant ($E_s = 10^4$ ou $\varepsilon_f = 0.4$) : il encaisse la
        vague et vibre. Plus fragile ($\varepsilon_f = 0.08$) : il se brise en plusieurs morceaux.
  \item Remplacez le bloc qui tombe par une rupture de barrage (colonne d'eau à gauche,
        \code{water_x0, water_x1 = 0.03, 0.30}, du sol au sommet) : la vague frappe le pilier et
        passe sous le tablier ; observez la poussée vers le haut.
  \item Un bâtiment : plusieurs poutres et poteaux (appelez \code{init_beam} plusieurs fois sur des
        plages d'indices, ou généralisez-le à une liste de rectangles).
\end{itemize}

% ================================================================
\section{Limites et pour aller plus loin}

\begin{enumerate}
  \item \textbf{Deux champs de vitesse.} Pour une interface nette et un glissement eau/solide, on
        garde deux grilles, $(m^f, v^f)$ et $(m^s, v^s)$, et l'on n'impose l'égalité des vitesses
        que dans la direction normale à l'interface, seulement si les deux phases se rapprochent
        (condition de contact). La normale se calcule comme le gradient de la masse solide sur la
        grille. C'est le schéma \og multi-vitesses \fg{} de Bardenhagen et al. et le premier pas vers
        un vrai frottement.
  \item \textbf{CPIC / fissures nettes.} Pour que deux fragments à moins d'une cellule ne se
        \og voient \fg{} plus, chaque n\oe ud porte deux vitesses (une par côté de la fissure) et les
        particules choisissent la leur selon une fonction de couleur.
  \item \textbf{Plasticité.} Entre l'élastique et la rupture, un matériau réel plastifie. Le même
        écrêtage des $s_k$ (borné à $[1-\theta_c,\,1+\theta_s]$ au lieu de $[0.1, 1]$) donne la
        plasticité de la neige ; Drucker-Prager donne le sable.
  \item \textbf{Réseau de liaisons.} Remplacer $D_p$ par des liaisons $B_{pq}$ entre particules
        voisines, cassées quand $\epsilon_{pq} > \epsilon_{\rm crit}$, rend la fissure discrète et
        orientée. Il faut une liste de voisins (fixée à l'initialisation) et une force de rappel
        $\mathbf f_{pq}$ le long de chaque liaison, en plus ou à la place de la contrainte MPM.
  \item \textbf{3D.} Comme pour le fluide : vecteurs à trois composantes, matrices $3\times3$,
        stencil $3\times3\times3$ ; \code{ti.svd} fonctionne en 3D et les formules ne changent pas.
\end{enumerate}

% ================================================================
\appendix
\section{Code complet : \texttt{MpM/mpm\_solid.py}}
\label{app:solid}

\begin{lstlisting}
# MpM/mpm_solid.py -- Solide MLS-MPM 2D : élasticité corotationnelle + endommagement + rupture
#
# Même grille et mêmes poids B-spline quadratiques que le fluide APIC (APIC/APIC.py).
# Chaque particule solide porte :
#   x, v, C   : position, vitesse, matrice affine APIC (comme le fluide)
#   F         : gradient de déformation (2x2), F = I au repos
#   D         : variable d'endommagement dans [0, 1]  (0 = intact, 1 = rompu)
#   broken    : 1 si la particule a perdu toute cohésion (rupture)

import taichi as ti


@ti.func
def kirchhoff_stress(F, mu: float, la: float):
    """tau = P F^T du modèle corotationnel : 2 mu (F - R) F^T + la J (J - 1) I."""
    U, sig, V = ti.svd(F)
    R = U @ V.transpose()
    J = F.determinant()
    I = ti.Matrix.identity(ti.f32, 2)
    return 2.0 * mu * (F - R) @ F.transpose() + la * J * (J - 1.0) * I


@ti.kernel
def P2G_solid(grid_m: ti.template(), grid_v: ti.template(),
              x: ti.template(), v: ti.template(), C: ti.template(), F: ti.template(),
              D: ti.template(), broken: ti.template(),
              inv_dx: float, dt: float, dx: float,
              mu: float, la: float, p_mass: float, p_vol: float):
    """Ne remet PAS la grille à zéro : on peut l'appeler après le P2G du fluide (grille partagée)."""
    for p in x:
        base = (x[p] * inv_dx - 0.5).cast(int)
        fx = x[p] * inv_dx - base
        w = [0.5 * (1.5 - fx)**2,
             0.75 - (fx - 1.0)**2,
             0.5 * (fx - 0.5)**2]

        k = 1.0 - D[p]
        if broken[p] == 1:
            k = 1.0
        tau = kirchhoff_stress(F[p], k * mu, k * la)

        affine = (-dt * p_vol * 4.0 * inv_dx * inv_dx) * tau + p_mass * C[p]

        for i, j in ti.static(ti.ndrange(3, 3)):
            offset = ti.Vector([i, j])
            d_pos = (offset - fx) * dx
            weight = w[i].x * w[j].y
            grid_v[base + offset] += weight * (p_mass * v[p] + affine @ d_pos)
            grid_m[base + offset] += weight * p_mass


@ti.kernel
def grid_update(grid_m: ti.template(), grid_v: ti.template(), mask: ti.template(),
                dt: float, g: float, bound: int, n_grid: int):
    for i, j in grid_m:
        if grid_m[i, j] > 0:
            grid_v[i, j] /= grid_m[i, j]
            grid_v[i, j].y -= dt * g
            if i < bound and grid_v[i, j].x < 0:
                grid_v[i, j].x = 0.0
            if i > n_grid - bound and grid_v[i, j].x > 0:
                grid_v[i, j].x = 0.0
            if j < bound and grid_v[i, j].y < 0:
                grid_v[i, j].y = 0.0
            if j > n_grid - bound and grid_v[i, j].y > 0:
                grid_v[i, j].y = 0.0
            if mask[i, j] == 1:
                grid_v[i, j] = [0.0, 0.0]


@ti.kernel
def G2P_solid(grid_v: ti.template(),
              x: ti.template(), v: ti.template(), C: ti.template(), F: ti.template(),
              D: ti.template(), broken: ti.template(),
              inv_dx: float, dt: float, dx: float, bound: int,
              eps0: float, epsf: float, use_damage: int, use_rupture: int):
    I = ti.Matrix.identity(ti.f32, 2)
    for p in x:
        base = (x[p] * inv_dx - 0.5).cast(int)
        fx = x[p] * inv_dx - base
        w = [0.5 * (1.5 - fx)**2,
             0.75 - (fx - 1.0)**2,
             0.5 * (fx - 0.5)**2]

        v_new = ti.Vector.zero(ti.f32, 2)
        C_new = ti.Matrix.zero(ti.f32, 2, 2)
        for i, j in ti.static(ti.ndrange(3, 3)):
            offset = ti.Vector([i, j])
            d_pos = (offset - fx) * dx
            weight = w[i].x * w[j].y
            g_v = grid_v[base + offset]
            v_new += weight * g_v
            C_new += weight * g_v.outer_product(d_pos) * (4.0 * inv_dx * inv_dx)
        v[p] = v_new
        C[p] = C_new

        F_new = (I + dt * C_new) @ F[p]

        if broken[p] == 1:
            U, sig, V = ti.svd(F_new)
            for d in ti.static(range(2)):
                sig[d, d] = ti.math.clamp(sig[d, d], 0.1, 1.0)
            F_new = U @ sig @ V.transpose()

        elif use_damage == 1:
            U, sig, V = ti.svd(F_new)
            eps = ti.max(sig[0, 0], sig[1, 1]) - 1.0
            D_new = ti.math.clamp((eps - eps0) / (epsf - eps0), 0.0, 1.0)
            D[p] = ti.max(D[p], D_new)
            if use_rupture == 1 and D[p] >= 1.0:
                broken[p] = 1
                F_new = I

        F[p] = F_new
        x[p] += dt * v[p]
        x[p] = ti.math.clamp(x[p], bound * dx, 1.0 - bound * dx)


@ti.kernel
def init_beam(x: ti.template(), v: ti.template(), C: ti.template(), F: ti.template(),
              D: ti.template(), broken: ti.template(),
              x0: float, y0: float, n_px: int, spacing: float):
    for p in x:
        i = p % n_px
        j = p // n_px
        x[p] = [x0 + (i + 0.5) * spacing, y0 + (j + 0.5) * spacing]
        v[p] = [0.0, 0.0]
        C[p] = ti.Matrix.zero(ti.f32, 2, 2)
        F[p] = ti.Matrix.identity(ti.f32, 2)
        D[p] = 0.0
        broken[p] = 0


@ti.kernel
def solid_colors(F: ti.template(), D: ti.template(), broken: ti.template(),
                 col: ti.template(), mode: int, eps_scale: float):
    for p in D:
        if mode == 0:
            if broken[p] == 1:
                col[p] = ti.Vector([0.90, 0.15, 0.15])
            else:
                col[p] = ti.Vector([0.75, 0.75, 0.75]) * (1 - D[p]) + ti.Vector([1.0, 0.85, 0.1]) * D[p]
        else:
            U, sig, V = ti.svd(F[p])
            eps = ti.max(sig[0, 0], sig[1, 1]) - 1.0
            t = ti.math.clamp(0.5 + 0.5 * eps / eps_scale, 0.0, 1.0)
            col[p] = ti.Vector([0.2, 0.4, 1.0]) * (1 - t) + ti.Vector([1.0, 0.25, 0.1]) * t


@ti.kernel
def solid_stats(D: ti.template(), broken: ti.template()) -> ti.types.vector(2, ti.f32):
    n_broken = 0
    D_max = 0.0
    for p in D:
        n_broken += broken[p]
        ti.atomic_max(D_max, D[p])
    return ti.Vector([float(n_broken), D_max])
\end{lstlisting}

% ================================================================
\section{Code complet : \texttt{MpM/demo\_solid.py}}
\label{app:demo}

\begin{lstlisting}
# MpM/demo_solid.py -- Partie I : une poutre MPM seule (sans fluide)
# Lancer depuis la racine du projet :  python -m MpM.demo_solid
# Touches : 1/2/3 modèle, ESPACE couleur, R reset, ESC quitter. Curseur "charge" = multiple de g.

import taichi as ti

try:
    from MpM.mpm_solid import *
except ImportError:          # lancé directement depuis le dossier MpM/
    from mpm_solid import *

ti.init(arch=ti.gpu)

n_grid = 128
dx = 1.0 / n_grid
inv_dx = float(n_grid)
bound = 3
render_substep = 20

rho_s = 2.0
E_s = 3000.0
nu_s = 0.3
mu_s = E_s / (2 * (1 + nu_s))
la_s = E_s * nu_s / ((1 + nu_s) * (1 - 2 * nu_s))
g = 9.81

eps0 = 0.05                          # début de l'endommagement (allongement principal)
epsf = 0.20                          # rupture complète (petit = fragile, grand = ductile)

c_p = ((la_s + 2 * mu_s) / rho_s) ** 0.5
dt = 0.4 * dx / c_p

p_spacing = 0.5 * dx                 # 2 x 2 particules par cellule
p_vol = p_spacing ** 2
p_mass = p_vol * rho_s

beam_x0, beam_x1 = 0.15, 0.85
beam_y0, beam_y1 = 0.55, 0.61
pillar_w = 0.06

n_px = int(round((beam_x1 - beam_x0) / p_spacing))
n_py = int(round((beam_y1 - beam_y0) / p_spacing))
n_solid = n_px * n_py

x_s = ti.Vector.field(2, ti.f32, n_solid)
v_s = ti.Vector.field(2, ti.f32, n_solid)
C_s = ti.Matrix.field(2, 2, ti.f32, n_solid)
F_s = ti.Matrix.field(2, 2, ti.f32, n_solid)
D_s = ti.field(ti.f32, n_solid)
broken_s = ti.field(ti.i32, n_solid)
col_s = ti.Vector.field(3, ti.f32, n_solid)

grid_v = ti.Vector.field(2, ti.f32, (n_grid, n_grid))
grid_m = ti.field(ti.f32, (n_grid, n_grid))
mask = ti.field(ti.i32, (n_grid, n_grid))
img = ti.Vector.field(3, ti.f32, (n_grid, n_grid))


@ti.kernel
def init_mask():
    for i, j in mask:
        pos = ti.Vector([i, j]) * dx
        left = beam_x0 <= pos.x <= beam_x0 + pillar_w
        right = beam_x1 - pillar_w <= pos.x <= beam_x1
        mask[i, j] = 1 if ((left or right) and pos.y <= beam_y1) else 0


@ti.kernel
def clear_grid():
    for i, j in grid_m:
        grid_v[i, j] = [0.0, 0.0]
        grid_m[i, j] = 0.0


@ti.kernel
def render_background():
    for i, j in img:
        img[i, j] = ti.Vector([0.35, 0.35, 0.35]) if mask[i, j] == 1 else ti.Vector([0.02, 0.02, 0.08])


def init_scene():
    init_beam(x_s, v_s, C_s, F_s, D_s, broken_s, beam_x0, beam_y0, n_px, p_spacing)
    init_mask()


def substep(load, use_damage, use_rupture):
    clear_grid()
    P2G_solid(grid_m, grid_v, x_s, v_s, C_s, F_s, D_s, broken_s,
              inv_dx, dt, dx, mu_s, la_s, p_mass, p_vol)
    grid_update(grid_m, grid_v, mask, dt, load * g, bound, n_grid)
    G2P_solid(grid_v, x_s, v_s, C_s, F_s, D_s, broken_s,
              inv_dx, dt, dx, bound, eps0, epsf, use_damage, use_rupture)


def main():
    window = ti.ui.Window("MPM 2D - poutre : élasticité, endommagement, rupture", res=(700, 700))
    canvas = window.get_canvas()
    gui = window.get_gui()
    model = 3
    color_mode = 0
    load = 1.0
    t_sim = 0.0
    init_scene()
    render_background()

    while window.running:
        while window.get_event(ti.ui.PRESS):
            k = window.event.key
            if k == ti.ui.ESCAPE:
                window.running = False
            elif k == ti.ui.SPACE:
                color_mode = 1 - color_mode
            elif k == 'r':
                init_scene()
                t_sim = 0.0
            elif k in ('1', '2', '3'):
                model = int(k)

        use_damage = 1 if model >= 2 else 0
        use_rupture = 1 if model == 3 else 0
        for _ in range(render_substep):
            substep(load, use_damage, use_rupture)
            t_sim += dt

        stats = solid_stats(D_s, broken_s)
        solid_colors(F_s, D_s, broken_s, col_s, color_mode, epsf)
        canvas.set_image(img)
        canvas.circles(x_s, radius=0.6 * p_spacing, per_vertex_color=col_s)

        gui.begin("MPM solide", 0.02, 0.02, 0.46, 0.30)
        load = gui.slider_float("charge (x g)", load, 0.0, 6.0)
        gui.text(f"modele : {['', 'elastique', 'endommagement', 'endommagement + rupture'][model]}   (touches 1/2/3)")
        gui.text(f"dt = {dt:.2e}   t = {t_sim:.3f} s")
        gui.text(f"D max = {stats[1]:.2f}   rompues = {int(stats[0])}/{n_solid}")
        gui.text("couleur : " + ("endommagement" if color_mode == 0 else "deformation principale") + "  (ESPACE)")
        gui.text("R : reset      ESC : quitter")
        gui.end()
        window.show()


if __name__ == "__main__":
    main()
\end{lstlisting}

% ================================================================
\section{Code complet : \texttt{main\_fsi.py}}
\label{app:fsi}

\begin{lstlisting}
# main_fsi.py -- Partie II : eau APIC + pont MPM sur une grille commune
# Touches : ESPACE couleur du solide, R reset, ESC quitter.

import taichi as ti
from APIC.APIC import P2G, G2P, grid_step
from MpM.mpm_solid import (P2G_solid, G2P_solid, init_beam, solid_colors, solid_stats)

ti.init(arch=ti.gpu)

n_grid = 128
dx = 1.0 / n_grid
inv_dx = float(n_grid)
bound = 3
render_substep = 20

# fluide (comme main.py)
rho_f = 1.0
E_f = 400.0
p_vol_f = (0.5 * dx) ** 2
p_mass_f = p_vol_f * rho_f

water_x0, water_x1, water_y0, water_y1 = 0.30, 0.70, 0.62, 0.95
pool_x0, pool_x1, pool_y0, pool_y1 = 0.03, 0.97, 0.03, 0.22
n_water = int(4 * (water_x1 - water_x0) * (water_y1 - water_y0) * n_grid * n_grid)
n_pool = int(4 * (pool_x1 - pool_x0) * (pool_y1 - pool_y0) * n_grid * n_grid)
n_fluid = n_water + n_pool

# solide (comme MpM/demo_solid.py)
rho_s = 2.0
E_s = 3000.0
nu_s = 0.3
mu_s = E_s / (2 * (1 + nu_s))
la_s = E_s * nu_s / ((1 + nu_s) * (1 - 2 * nu_s))
eps0, epsf = 0.05, 0.20
use_damage, use_rupture = 1, 1

p_spacing = 0.5 * dx
p_vol_s = p_spacing ** 2
p_mass_s = p_vol_s * rho_s

beam_x0, beam_x1 = 0.15, 0.85
beam_y0, beam_y1 = 0.40, 0.46
pillar_w = 0.06
n_px = int(round((beam_x1 - beam_x0) / p_spacing))
n_py = int(round((beam_y1 - beam_y0) / p_spacing))
n_solid = n_px * n_py

# pas de temps commun
g = 9.81
c_f = (E_f / rho_f) ** 0.5
c_p = ((la_s + 2 * mu_s) / rho_s) ** 0.5
dt = 0.4 * dx / max(c_f, c_p)

# champs fluide
x_f = ti.Vector.field(2, ti.f32, n_fluid)
v_f = ti.Vector.field(2, ti.f32, n_fluid)
C_f = ti.Matrix.field(2, 2, ti.f32, n_fluid)
J_f = ti.field(ti.f32, n_fluid)

# champs solide
x_s = ti.Vector.field(2, ti.f32, n_solid)
v_s = ti.Vector.field(2, ti.f32, n_solid)
C_s = ti.Matrix.field(2, 2, ti.f32, n_solid)
F_s = ti.Matrix.field(2, 2, ti.f32, n_solid)
D_s = ti.field(ti.f32, n_solid)
broken_s = ti.field(ti.i32, n_solid)
col_s = ti.Vector.field(3, ti.f32, n_solid)

# grille commune
grid_v = ti.Vector.field(2, ti.f32, (n_grid, n_grid))
grid_m = ti.field(ti.f32, (n_grid, n_grid))
solid = ti.field(ti.f32, (n_grid, n_grid))            # 1 = pilier
img = ti.Vector.field(3, ti.f32, (n_grid, n_grid))


@ti.kernel
def init_fluid():
    for p in x_f:
        if p < n_water:
            x_f[p] = [water_x0 + ti.random() * (water_x1 - water_x0),
                      water_y0 + ti.random() * (water_y1 - water_y0)]
        else:
            x_f[p] = [pool_x0 + ti.random() * (pool_x1 - pool_x0),
                      pool_y0 + ti.random() * (pool_y1 - pool_y0)]
        v_f[p] = [0.0, 0.0]
        C_f[p] = ti.Matrix.zero(ti.f32, 2, 2)
        J_f[p] = 1.0
    for i, j in solid:
        pos = ti.Vector([i, j]) * dx
        left = beam_x0 <= pos.x <= beam_x0 + pillar_w
        right = beam_x1 - pillar_w <= pos.x <= beam_x1
        solid[i, j] = 1.0 if ((left or right) and pos.y <= beam_y1) else 0.0


@ti.kernel
def advect_fluid():
    for p in x_f:
        x_f[p] += dt * v_f[p]
        x_f[p] = ti.math.clamp(x_f[p], bound * dx, 1.0 - bound * dx)


@ti.kernel
def render_background():
    for i, j in img:
        img[i, j] = ti.Vector([0.35, 0.35, 0.35]) if solid[i, j] == 1 else ti.Vector([0.02, 0.02, 0.08])


def init_scene():
    init_fluid()
    init_beam(x_s, v_s, C_s, F_s, D_s, broken_s, beam_x0, beam_y0, n_px, p_spacing)


def substep():
    P2G(grid_m, grid_v, x_f, v_f, C_f, J_f, inv_dx, dt, dx, E_f, p_mass_f, p_vol_f)      # 1
    P2G_solid(grid_m, grid_v, x_s, v_s, C_s, F_s, D_s, broken_s,
              inv_dx, dt, dx, mu_s, la_s, p_mass_s, p_vol_s)                              # 2
    grid_step(grid_m, grid_v, solid, dt, g, bound, n_grid)                                # 3
    G2P(grid_m, grid_v, x_f, v_f, C_f, J_f, inv_dx, dt, dx, E_f, p_mass_f, p_vol_f)      # 4
    G2P_solid(grid_v, x_s, v_s, C_s, F_s, D_s, broken_s,
              inv_dx, dt, dx, bound, eps0, epsf, use_damage, use_rupture)                 # 5
    advect_fluid()                                                                        # 6


def main():
    window = ti.ui.Window("APIC + MPM : eau et pont", res=(700, 700))
    canvas = window.get_canvas()
    gui = window.get_gui()
    color_mode = 0
    t_sim = 0.0
    init_scene()
    render_background()

    while window.running:
        while window.get_event(ti.ui.PRESS):
            k = window.event.key
            if k == ti.ui.ESCAPE:
                window.running = False
            elif k == ti.ui.SPACE:
                color_mode = 1 - color_mode
            elif k == 'r':
                init_scene()
                t_sim = 0.0

        for _ in range(render_substep):
            substep()
            t_sim += dt

        stats = solid_stats(D_s, broken_s)
        solid_colors(F_s, D_s, broken_s, col_s, color_mode, epsf)
        canvas.set_image(img)
        canvas.circles(x_f, radius=0.5 * p_spacing, color=(0.35, 0.65, 1.0))
        canvas.circles(x_s, radius=0.6 * p_spacing, per_vertex_color=col_s)

        gui.begin("APIC + MPM", 0.02, 0.02, 0.46, 0.22)
        gui.text(f"dt = {dt:.2e}   t = {t_sim:.3f} s")
        gui.text(f"D max = {stats[1]:.2f}   rompues = {int(stats[0])}/{n_solid}")
        gui.text("couleur : " + ("endommagement" if color_mode == 0 else "deformation principale") + "  (ESPACE)")
        gui.text("R : reset      ESC : quitter")
        gui.end()
        window.show()


if __name__ == "__main__":
    main()
\end{lstlisting}

\end{document}
